\section*{Hoofdstuk \ref{ch:b1:atmospheric-circulation} --- De algemene circulatie van de atmosfeer}
% values: calc/ch28.py

\begin{solution}{exo:b1:atmospheric-circulation:1}
(a)~$f = 2\Omega\sin\varphi$: $0$ aan de evenaar, \qty{7.29e-5}{s^{-1}} op \ang{30},
\qty{1.031e-4}{s^{-1}} op \ang{45} en \qty{1.458e-4}{s^{-1}} op de pool.
(b)~$fV = \num{1.031e-4} \times 10 = \qty{1.03e-3}{m.s^{-2}}$, ongeveer \num{1.05e-4}
van $g$. (c)~Op het noordelijk halfrond wijkt de afbuiging naar rechts af van de beweging: een wind die naar het noorden waait, wordt naar het oosten geduwd.
\end{solution}

\begin{solution}{exo:b1:atmospheric-circulation:2}
(a)~$|\nabla p|/\rho = (400/\num{4e5})/1.225 = \qty{8.16e-4}{m.s^{-2}}$.
(b)~$10/\num{8.16e-4} = \qty{12250}{s}$, \qty{3.4}{h}. (c)~De Coriolisversnelling
van een wind van \qty{10}{m/s} op \ang{45} N is \qty{1.03e-3}{m.s^{-2}}: de twee zijn even groot
(verhouding $0.79$). Binnen enkele uren nadat de lucht begint te bewegen, voelt zij een Coriolis-kracht die de drukgradiënt in evenwicht kan houden: de wind stelt zich in op ongeveer de geostrofische snelheid, hier $\num{8.16e-4}/\num{1.031e-4} = \qty{7.9}{m/s}$, in plaats van langs de gradiënt te blijven versnellen.
\end{solution}

\begin{solution}{exo:b1:atmospheric-circulation:3}
(a)~$f = 2 \times \num{7.292e-5} \times \sin\ang{40} = \qty{9.374e-5}{s^{-1}}$;
$|\vec v_g| = (400/\num{2.5e5})/(1.225 \times \num{9.374e-5}) = \qty{13.9}{m/s}$, of
$13.93/0.5144 = \qty{27.1}{kn}$. (b)~Op het noordelijk halfrond houdt de wind de lage druk aan zijn linkerhand: met de lage druk ten zuiden waait hij naar het westen, dus uit het oosten
(\omterm{def:b1:spherical-astronomy:horizontal}{azimut} \ang{90}). (c)~De staf wijst vanaf de stationscirkel naar het oosten, de wind tegemoet; \qty{27.1}{kn} wordt afgerond tot \qty{25}{kn}, twee volle veren en een halve veer, getekend aan de met de klok mee gelegen zijde van de staf, hier de zuidzijde, naar de lage druk toe.
\end{solution}

\begin{solution}{exo:b1:atmospheric-circulation:4}
(a)~\ang{15} N: noordoostelijke \omterm{def:b1:atmospheric-circulation:winds}{passaten}; \ang{15} Z: zuidoostelijke \omterm{def:b1:atmospheric-circulation:winds}{passaten};
\ang{45} N en \ang{45} Z: \omterm{def:b1:atmospheric-circulation:winds}{westenwinden}. (b)~Op \ang{15} keert de oppervlaktelucht van de \omterm{def:b1:atmospheric-circulation:cells}{Hadley-cel} terug naar de evenaar en wordt door de Coriolis-kracht naar het westen gedraaid
(rechts op het noordelijk halfrond, links op het zuidelijk): oostelijke wind, negatief $u$. Op \ang{45} beweegt de oppervlaktelucht van de \omterm{def:b1:atmospheric-circulation:cells}{Ferrel-cel} naar de pool en wordt naar het oosten gedraaid: westelijke wind, positief $u$. (c)~Van de drie breedtegraden is de druk het hoogst op \ang{30} N,
\qty{1016.8}{hPa}: de subtropische hogedrukgordel ligt nabij \ang{30} N (het volledige profiel plaatst haar maximum nabij \ang{35} N).
\end{solution}

\begin{solution}{exo:b1:atmospheric-circulation:5}
(a)~$|\nabla p|/\rho = (200/\num{5e5})/1.225 = \qty{3.27e-4}{m.s^{-2}}$, gedeeld door
$f$: \qty{2.6}{m/s} op \ang{60}, \qty{2.9}{m/s} op \ang{50}, \qty{4.5}{m/s} op
\ang{30}, \qty{6.5}{m/s} op \ang{20}, \qty{12.9}{m/s} op \ang{10} en
\qty{25.7}{m/s} op \ang{5}. (b)~De Coriolis-kracht draait een wind in een tijd van de orde van
$1/f$: \qty{2.2}{h} op \ang{60}, \qty{11}{h} op \ang{10}, \qty{22}{h} op \ang{5}, en
oneindig lang aan de evenaar. Nabij de evenaar verandert het drukveld, en werkt wrijving, lang voordat de wind langs de isobaren kan worden gedraaid: het evenwicht komt nooit tot stand, en de lucht stroomt deels langs de drukgradiënt, naar lagere druk, met bescheiden snelheden.
\end{solution}

\begin{solution}{exo:b1:atmospheric-circulation:6}
(a)~$u = 464.6\,\sin^2\varphi/\cos\varphi$: \qty{14.2}{m/s} op \ang{10},
\qty{32.2}{m/s} op \ang{15}, \qty{57.8}{m/s} op \ang{20}, \qty{134.1}{m/s} op
\ang{30}. (b)~$u = 464.6\,(\cos^2\ang{10} - \cos^2\varphi)/\cos\varphi$:
\qty{17.7}{m/s} op \ang{15}, \qty{42.9}{m/s} op \ang{20}, \qty{117.9}{m/s} op
\ang{30}. (c)~Lucht die verder van de evenaar begint, begint dichter bij de as, met minder impulsmoment over: zij komt met een zwakkere westenwind aan op elke \omterm{def:b1:spherical-astronomy:geographic}{breedtegraad}. Het uitspreiden van de opstijgende tak verlaagt de voorspelde \omterm{def:b1:atmospheric-circulation:jet}{straalstroom}, maar niet genoeg om de waargenomen \qty{40}{m/s} te bereiken.
\end{solution}

\begin{solution}{exo:b1:atmospheric-circulation:7}
(a)~De wind waait naar het noordoosten. Met de rug naar de wind, dus gericht op het noordoosten, ligt de lage druk links: naar het noordwesten. (b)~De wind waait naar het zuidoosten; op het zuidelijk halfrond ligt de lage druk rechts van iemand met de rug naar de wind: naar het zuidwesten. (c)~De oppervlaktewind wordt van de isobaren naar de lage druk toe gedraaid, dat wil zeggen linksom ten opzichte van de \omterm{def:b1:atmospheric-circulation:geostrophic}{geostrofische wind} op het noordelijk halfrond. De \omterm{def:b1:atmospheric-circulation:geostrophic}{geostrofische wind}, en de lage druk die er rechthoekig links van ligt, zijn daarom rechtsom gedraaid ten opzichte van het antwoord van (a): de lage druk ligt tussen het noordwesten en het noorden, op de \omterm{def:b1:spherical-astronomy:horizontal}{azimut} $\ang{315} + \alpha$ voor een kruishoek $\alpha$.
\end{solution}

\begin{solution}{exo:b1:atmospheric-circulation:8}
(a)~Het oppervlak tussen de evenaar en \ang{30} N is $2\pi a_\oplus^2\sin\ang{30} =
\qty{1.275e14}{m^2}$; de massa boven een vierkante meter tussen de twee niveaus is
$\Delta p/g = \num{9e4}/9.81 = \qty{9174}{kg}$; het totaal is
\qty{1.170e18}{kg}. (b)~$\num{1.170e18}/\num{8.7e10} = \qty{1.34e7}{s}$, ongeveer
\num{156} dagen, vijf maanden. (c)~$\num{1.170e18}/\num{21.4e10} = \qty{5.47e6}{s}$,
ongeveer \num{63} dagen, twee maanden.
\end{solution}

\begin{solution}{exo:b1:atmospheric-circulation:9}
(a)~$\num{0.01} \times \num{86400} = \qty{864}{m}$ per dag. (b)~$0.864 \times 9.8 =
\qty{8.5}{K}$ per dag. (c)~De lucht koelt af door infraroodstraling uit te zenden: de atmosfeer verliest meer energie door straling dan zij absorbeert, en de daling van de lucht, die haar door compressie opwarmt, vult het verlies aan. (d)~Dalende lucht warmt op en wordt relatief droger: haar wolken verdampen en regen kan niet ontstaan. De \omterm{def:b1:atmospheric-circulation:cells}{subtropische hogedrukgordels} zijn gordels van heldere hemel en weinig regen, waar op het land hete woestijnen ontstaan.
\end{solution}

\begin{solution}{exo:b1:atmospheric-circulation:10}
(a)~$V^2 + frV - fr|\vec v_g| = 0$ met $fr = \qty{55.85}{m/s}$:
$V = \tfrac12\bigl(-55.85 + \sqrt{55.85^2 + 4 \times 55.85 \times 15}\bigr) =
\qty{12.3}{m/s}$, \qty{82}{\percent} van de \omterm{def:b1:atmospheric-circulation:geostrophic}{geostrofische wind}. (b)~$V^2 - frV +
fr|\vec v_g| = 0$ heeft reële wortels alleen als $(fr)^2 \geq 4fr|\vec v_g|$, dat wil zeggen
$|\vec v_g| \leq fr/4 = \qty{14.0}{m/s}$. (c)~Nabij het centrum van een hoge druk, waar $r$ klein is, is de grens $fr/4$ klein: de drukgradiënt moet zwak zijn, de isobaren ver uit elkaar en de winden licht. Een depressie kent geen zo'n grens: haar isobaren kunnen dicht bij het centrum samendrommen, met sterke winden.
\end{solution}

\begin{solution}{exo:b1:atmospheric-circulation:11}
(a)~Voor kleine $\varphi$ is $\sin\varphi \approx \varphi$ en $\cos\varphi \approx 1$,
zodat $u \approx \Omega a_\oplus\varphi^2$. (b)~$\ang{10} = \qty{0.1745}{rad}$:
$464.6 \times 0.1745^2 = \qty{14.15}{m/s}$, tegen \qty{14.23}{m/s} exact,
\qty{0.5}{\percent} minder. (c)~Het numeriek oplossen van $464.6\sin^2\varphi/\cos\varphi = 27.0$
geeft $\varphi = \ang{13.7}$. (d)~De ring haalt de werkelijke jaargemiddelde \omterm{def:b1:atmospheric-circulation:jet}{straalstroom} al op minder dan \ang{14}: behoud van impulsmoment beschrijft de bovenste tak van de \omterm{def:b1:atmospheric-circulation:cells}{Hadley-cel} alleen in de diepe tropen; poolwaarts van ongeveer \ang{15} nemen de wervels het over.
\end{solution}

\begin{solution}{exo:b1:atmospheric-circulation:12}
(a)~Gemiddelde $31.7/12 = \ang{2.6}$ N; spreiding van \ang{12.8} Z tot \ang{10.9} N, dat wil zeggen
\ang{23.7}. (b)~Het noordelijkste punt komt in september, drie maanden na de zomerzonnewende van juni; het zuidelijkste in februari, twee maanden na de winterzonnewende van december. De vertraging komt door warmteopslag, vooral in de \omterm{def:b1:fluid-envelopes-structure:ocean-layers}{gemengde laag} van de oceaan (\cref{ch:b1:fluid-envelopes-structure}), die langzaam opwarmt en afkoelt. (c)~Het zonaal gemiddelde middelt alle lengtegraden: als boven de continenten de regenband ver het zomerhalfrond in zwaait, kan het gemiddelde over alle lengtegraden de evenaar kruisen, ook al doet de band boven één oceaan dat niet.
\end{solution}

\begin{solution}{pb:b1:atmospheric-circulation:1}
\begin{enumerate}
  \item $\rho = p/(R_dT) = \num{100600}/(287.05 \times 283.15) =
        \qty{1.238}{kg.m^{-3}}$.
  \item $|\nabla p| = \qty{400}{Pa}/\qty{2.5e5}{m} = \qty{1.6e-3}{Pa.m^{-1}}$.
  \item $|\nabla p|/\rho = \num{1.6e-3}/1.238 = \qty{1.293e-3}{m.s^{-2}}$.
  \item $\Omega = 2\pi/\num{86164.2} = \qty{7.292e-5}{rad.s^{-1}}$;
        $f = 2\Omega\sin\ang{50} = \qty{1.117e-4}{s^{-1}}$.
  \item $|\vec v_g| = \num{1.293e-3}/\num{1.117e-4} = \qty{11.57}{m/s}$.
  \item $11.57/0.514444 = \qty{22.5}{kn}$; $11.57 \times 3.6 = \qty{41.7}{km/h}$.
  \item De tijd is $|\vec v_g|/(|\nabla p|/\rho)$ en $|\vec v_g| = (|\nabla
        p|/\rho)/f$, dus is zij $1/f = \qty{8951}{s}$, \qty{2.49}{h}; en $2\pi/f =
        \qty{15.6}{h}$. Patronen die een dag of langer duren liggen daarom dicht bij het evenwicht.
  \item De \omterm{def:b1:atmospheric-circulation:pgf}{drukgradiëntkracht} wijst naar het noorden, naar de lage druk toe. De \omterm{def:b1:atmospheric-circulation:geostrophic}{geostrofische
        wind} waait langs de isobaren met de lage druk aan zijn linkerhand: naar het oosten,
        uit het westen, \omterm{def:b1:spherical-astronomy:horizontal}{azimut} \ang{270}: een westenwind.
  \item Met de rug naar een westenwind, gericht naar het oosten, ligt het noorden links,
        waar de lage druk ligt: de wet van Buys Ballot is vervuld.
  \item De staf wijst vanaf de stationscirkel naar het westen, naar de kant waar de wind
        vandaan komt; \qty{22.5}{kn} wordt afgerond tot \qty{20}{kn}, twee volle veren, die
        aan de met de klok mee gelegen zijde van de staf staan, de noordzijde, naar de lage druk toe.
  \item Op \ang{50} Z is $f < 0$: de wind houdt de lage druk aan zijn rechterhand, dus waait
        hij naar het westen, uit het oosten.
  \item $f(\ang{30}) = \qty{7.292e-5}{s^{-1}}$: \qty{17.7}{m/s}; $f(\ang{60}) =
        \qty{1.263e-4}{s^{-1}}$: \qty{10.2}{m/s}.
  \item $f(\ang{5}) = \qty{1.271e-5}{s^{-1}}$ geeft \qty{101.7}{m/s}, een wind die bij zulke gradiënten nooit is waargenomen. De draaitijd $1/f$ is \qty{22}{h}:
        het evenwicht heeft geen tijd om tot stand te komen, en de formule geldt niet.
  \item $9.0/11.57 = 0.778$; langs de isobaren $9.0\cos\ang{25} = \qty{8.16}{m/s}$,
        dwars op de isobaren naar de lage druk toe $9.0\sin\ang{25} = \qty{3.80}{m/s}$.
  \item De \omterm{def:b1:atmospheric-circulation:geostrophic}{geostrofische wind} waait naar \omterm{def:b1:spherical-astronomy:horizontal}{azimut} \ang{90}; gedraaid \ang{25} naar de lage druk in het noorden, waait de oppervlaktewind naar \ang{65},
        uit \ang{245}, het west-zuidwesten.
  \item $\tan\alpha = 1/(f\tau)$ geeft $\tau = 1/(\num{1.117e-4}\tan\ang{25}) =
        \qty{1.92e4}{s}$, \qty{5.3}{h}; het model voorspelt $11.57\cos\ang{25} =
        \qty{10.5}{m/s}$, meer dan de waargenomen \qty{9.0}{m/s}: echte wrijving is niet precies lineair in de wind en vertraagt de wind bij deze hoek meer dan het model.
  \item Door een cilinder met straal $R = \qty{500}{km}$ en hoogte $h =
        \qty{1}{km}$: $\dot m = 2\pi R h\rho u_r = 2\pi \times \num{5e5} \times 1000
        \times 1.238 \times 3.80 = \qty{1.48e10}{kg/s}$.
  \item $w = \dot m/(\rho\pi R^2) = 2hu_r/R = 2 \times 1000 \times 3.80/\num{5e5} =
        \qty{0.0152}{m/s}$, ongeveer \qty{1310}{m} per dag.
  \item $1.31 \times 9.8 = \qty{12.9}{K}$ per dag: de lucht bereikt al snel verzadiging,
        en de depressie vult zich met wolken en regen. Onder een \omterm{def:b1:atmospheric-circulation:cyclone}{anticycloon} stroomt de oppervlaktelucht naar buiten, daalt de lucht erboven en warmt zij op, en klaart de hemel op.
  \item $\Omega a_\oplus = \num{7.292e-5} \times \num{6.371e6} = \qty{464.6}{m/s}$
        op de evenaar; $464.6\cos\ang{30} = \qty{402.3}{m/s}$ op \ang{30} N.
  \item $u = 464.6\,\sin^2\varphi/\cos\varphi$: \qty{57.8}{m/s} op \ang{20} N en
        \qty{134.1}{m/s} op \ang{30} N.
  \item $u = 464.6\,(\cos^2\ang{10} - \cos^2\ang{30})/\cos\ang{30} =
        \qty{117.9}{m/s}$.
  \item $134.1/42.3 = 3.17$. De wervels van de middelbreedten, depressies en
        anticyclonen, nemen impulsmoment uit de bovenste tak en voeren het elders heen; en de lucht van de cel stijgt boven een band van breedtegraden in plaats van in rust op de evenaar. Het zonaal gemiddelde middelt ook een meanderende \omterm{def:b1:atmospheric-circulation:jet}{straalstroom} over alle lengtegraden.
  \item $M/M_0 = (402.3 + 42.3)\cos\ang{30}/464.6 = 0.829$: de lucht heeft
        \qty{17}{\percent} van het impulsmoment verloren dat zij als vrije ring had behouden.
  \item De \omterm{def:b1:atmospheric-circulation:geostrophic}{geostrofische wind} op P is \qty{11.6}{m/s}, uit het westen; de aan de grond gemeten wind is \qty{9.0}{m/s}, ongeveer vier vijfde daarvan, en door wrijving \ang{25} naar de lage druk toe gedraaid.
\end{enumerate}
\end{solution}
